\subsection{From a given power to the pulse area}

The pulse area is the only quantity the excitation depends on, and it follows
from the optical power in three steps: fix the scale of the intensity, convert
intensity into a Rabi frequency, integrate over the pulse. This section keeps
the three normalisations apart, because only one of them is physics.

\paragraph{What is normalised, and what is not.}
\begin{enumerate}
\item \emph{The mode.} $u(\mathbf 0)=1$, so the simulated intensity
      $\tilde I(\mathbf r,t)=|E|^2$ is in arbitrary units. This choice is free
      and drops out of every physical result.
\item \emph{The power.} The optical power $P$ in the profile is the single
      physical input. It fixes the scale factor $s$ in
      $I(\mathbf r,t)=s\,\tilde I(\mathbf r,t)$.
\item \emph{The reference.} A reference intensity $I_\mathrm{ref}$, an
      effective area $A_\mathrm{eff}$ and a nominal Rabi frequency
      $f_\mathrm{Rabi}$ are bookkeeping. They make the numbers readable and
      cancel from the pulse area.
\end{enumerate}

\paragraph{Step 1: the scale from the power.}
The power carried by the pattern is the area integral of the time average,
\begin{equation}
P=\int\!\langle I(\mathbf r)\rangle_t\,\mathrm d^2r
 =s\int\!\langle\tilde I(\mathbf r)\rangle_t\,\mathrm d^2r
 \equiv s\,\tilde P_\mathrm{prof},
\qquad\Longrightarrow\qquad
s=\frac{P}{\tilde P_\mathrm{prof}} .
\end{equation}
$\tilde P_\mathrm{prof}$ is known in closed form, and it is worth doing the
three integrals explicitly.

\subparagraph{(a) The power of one spot.}
The notation is the one of the previous section: $g_s=A_s\,u(\mathbf r-\mathbf
c_s)$ with the field amplitude $A_s=\sqrt{a_{x,i}a_{y,j}}$, and it is
convenient to abbreviate the intensity weight of a spot by
\begin{equation}
a_s\equiv A_s^2=a_{x,i}\,a_{y,j} .
\end{equation}
A single spot then carries $\int g_s^2\,\mathrm d^2r=a_sP_1$, with
$P_1=\int|u|^2\mathrm d^2r$.

\emph{Notation.} A bold symbol is a vector in the focal plane, the same letter
in italics its length: $r=|\mathbf r|$, $q=|\mathbf q|$, $d=|\mathbf d|$. All
integrals $\int\!\ldots\mathrm d^2r$ and $\int\!\ldots\mathrm d^2q$ run over the
plane, i.e.\ over the vector. The mode is radially symmetric, so $u$ is written
as a function of the \emph{scalar} radius, $u(r)\equiv u(|\mathbf r|)$, and the
same holds for $\tilde u(q)$ and for the overlap $O(d)$. For the Gaussian, substituting
$x=2r^2/w_0^2$,
\begin{equation}
P_1^\mathrm{G}=\int_0^\infty\! e^{-2r^2/w_0^2}\,2\pi r\,\mathrm dr
=\frac{\pi w_0^2}{2}\int_0^\infty\! e^{-x}\,\mathrm dx=\frac{\pi w_0^2}{2}.
\end{equation}
For the Airy mode, with $x=kr$ and the standard integral
$\int_0^\infty J_\nu^2(x)\,\mathrm dx/x=1/(2\nu)$,
\begin{gather}
P_1^\mathrm{A}=\int_0^\infty\!\left(\frac{2J_1(kr)}{kr}\right)^{\!2}2\pi r\,\mathrm dr
=\frac{8\pi}{k^2}\int_0^\infty\!\frac{J_1^2(x)}{x}\,\mathrm dx
=\frac{8\pi}{k^2}\cdot\frac12=\frac{4\pi}{k^2},\\
k=\frac{3.8317}{1.4830\,w_0}. \notag
\end{gather}
The Airy spot carries $4\pi/k^2\big/(\pi w_0^2/2)=1.20$ times the power of a
Gaussian of the same $1/e^2$ width, because of its rings.

\subparagraph{(b) Which cross terms survive the time average.}
The time average of the coherent sum was derived above,
\begin{equation}
\langle\tilde I(\mathbf r)\rangle_t=\sum_s g_s^2(\mathbf r)
+2\!\!\sum_{\substack{s<s'\\ k_s=k_{s'}}}\!\! g_s(\mathbf r)g_{s'}(\mathbf r)
\cos(\phi_s-\phi_{s'}) ,
\end{equation}
because a pair with $k_s\neq k_{s'}$ oscillates at $(k_s-k_{s'})f_0$ and
averages to zero over $T_0$, while a degenerate pair does not oscillate at all.
Integrating term by term over the plane gives
\begin{equation}
\tilde P_\mathrm{prof}=\sum_s a_s\,P_1
+2\!\!\sum_{\substack{s<s'\\ k_s=k_{s'}}}\!\!A_sA_{s'}\,
\cos(\phi_s-\phi_{s'})\,
\underbrace{\int\! u(\mathbf r-\mathbf c_s)u(\mathbf r-\mathbf c_{s'})\,
\mathrm d^2r}_{\textstyle O(d)} ,
\end{equation}
with $d=|\mathbf c_s-\mathbf c_{s'}|$. The first sum therefore runs over the
intensity weights $a_s=A_s^2$ and not over $a_s^2$, simply because
$g_s^2=A_s^2|u|^2$.

\subparagraph{(c) The overlap integral.}
For the Gaussian, completing the square with the midpoint
$\mathbf m=\tfrac12(\mathbf c_s+\mathbf c_{s'})$,
\begin{gather}
e^{-\frac{(\mathbf r-\mathbf c_s)^2}{w_0^2}}
e^{-\frac{(\mathbf r-\mathbf c_{s'})^2}{w_0^2}}
=e^{-\frac{d^2}{2w_0^2}}\;e^{-\frac{2(\mathbf r-\mathbf m)^2}{w_0^2}},\\
\Longrightarrow\quad
O^\mathrm{G}(d)=\frac{\pi w_0^2}{2}\,e^{-d^2/2w_0^2}
=P_1^\mathrm{G}\,e^{-d^2/2w_0^2}.
\end{gather}
For the Airy mode the direct integral is awkward, but Fourier space makes it a
one-liner. Three ingredients are needed.

\emph{(i) The mode is band limited.} The focal field is the Fourier transform
of the pupil field, and a uniformly illuminated circular pupil is a disc. Its
transform is
\begin{equation}
\frac{1}{(2\pi)^2}\int_{|\mathbf q|<k}\! e^{i\mathbf q\cdot\mathbf r}\,\mathrm d^2q
=\frac{k^2}{4\pi}\,\frac{2J_1(kr)}{kr}
=\frac{k^2}{4\pi}\,u(r),\qquad r=|\mathbf r| ,
\end{equation}
where the left side is an integral over the vector $\mathbf q$ while the right
side depends only on the length $r$, as it must for a radially symmetric
pupil. Read backwards,
\begin{equation}
\tilde u(q)=\int\! u(|\mathbf r|)\,e^{-i\mathbf q\cdot\mathbf r}\,\mathrm d^2r
=\frac{4\pi}{k^2}\,\Theta\!\left(k-q\right),
\qquad q=|\mathbf q| ,
\end{equation}
a \emph{flat} spectrum inside the cutoff $k=2\pi\,\mathrm{NA}/\lambda$ and zero
outside. The prefactor follows from $\mathbf q=0$, where the exponential is unity and
the plane integral becomes a radial one,
$\tilde u(0)=\int u\,\mathrm d^2r=2\pi\int_0^\infty u(r)\,r\,\mathrm dr
=\frac{4\pi}{k}\int_0^\infty J_1(kr)\,\mathrm dr=\frac{4\pi}{k^2}$,
using $\int_0^\infty J_1(x)\,\mathrm dx=1$.

\emph{(ii) The integral over a disc.} The angular integral gives a Bessel
function, $\int_0^{2\pi}e^{ix\cos\varphi}\,\mathrm d\varphi=2\pi J_0(x)$, and
the radial one closes with $\int_0^XJ_0(t)\,t\,\mathrm dt=XJ_1(X)$:
\begin{equation}
\int_{|\mathbf q|<k}\! e^{i\mathbf q\cdot\mathbf d}\,\mathrm d^2q
=\int_0^k\!\!\int_0^{2\pi}\! e^{i\,qd\cos\varphi}\,q\,\mathrm d\varphi\,\mathrm dq
=2\pi\!\int_0^k\! J_0(qd)\,q\,\mathrm dq
=\frac{2\pi k\,J_1(kd)}{d} .
\end{equation}
In the first step $\mathbf q\cdot\mathbf d=qd\cos\varphi$ was used, with
$\varphi$ the angle between the two vectors; from there on only the lengths
$q$ and $d$ appear.

\emph{(iii) The correlation theorem.} The overlap depends only on the spot
separation $\mathbf d=\mathbf c_{s'}-\mathbf c_s$, and since $u$ is real and
even the correlation equals the convolution,
\begin{equation}
O^\mathrm{A}(d)=\int\! u(|\mathbf r|)\,u(|\mathbf r-\mathbf d|)\,\mathrm d^2r
=\frac{1}{(2\pi)^2}\int\!\left|\tilde u(q)\right|^2
e^{i\mathbf q\cdot\mathbf d}\,\mathrm d^2q .
\end{equation}
The shift $\mathbf r-\mathbf d$ is a vector operation, but the result again
depends only on the distance $d=|\mathbf d|$.
Here comes the point: the spectrum is an indicator function, so squaring it
changes nothing, $\Theta^2=\Theta$. The autocorrelation therefore has the same
spectrum as the mode itself, only with a different prefactor, and putting
(i)--(iii) together,
\begin{equation}
O^\mathrm{A}(d)
=\frac{1}{(2\pi)^2}\left(\frac{4\pi}{k^2}\right)^{\!2}\frac{2\pi k J_1(kd)}{d}
=\frac{8\pi\,J_1(kd)}{k^3 d}
=\frac{4\pi}{k^2}\cdot\frac{2J_1(kd)}{kd}
=P_1^\mathrm{A}\,u(d) .
\end{equation}
The overlap of two Airy spots is the single-spot power times the Airy mode
evaluated at their distance -- \emph{including its sign}, which is negative
between the first and the second ring. Two checks: at $d=0$ it reduces to
$P_1^\mathrm{A}$ (Parseval), and against a numerical FFT it agrees to
$0.4\,\%$.

The same result read in the pupil: by Parseval the overlap in the focal plane
equals the overlap in the pupil plane. There the two spots are two plane waves
tilted against each other, which is a linear phase $e^{i\mathbf q\cdot\mathbf
d}$ across the pupil coordinate $\mathbf q$. Averaging that phase over the disc
is exactly integral (ii) and gives the visibility factor $2J_1(kd)/(kd)$: the
further apart the two spots, the more fringes fit across the pupil and the more
they cancel. In contrast to the Gaussian the decay is slow, $\propto
d^{-3/2}$, and the sign oscillates.

For $3\times4$ at the working point the degenerate corner pair sits
$d=3.305\,\mu$m apart, $u(d)=+0.063$, and with $A_sA_{s'}=a_s=1.1252$ it
contributes $2\cdot1.1252\cdot0.063=0.142$ against $\sum_sa_s=12.70$, i.e.
$+1.1\,\%$ at all phases zero.

The integral is never taken from the computation grid, which truncates about
$5\,\%$ of the Airy power at the usual margin.

\paragraph{Step 2: intensity to Rabi frequency.}
The two Raman legs, of intensities $\beta I$ and $(1-\beta)I$, drive
$|F{=}2,m\rangle\leftrightarrow|F{=}3,m\rangle$ via the $5P$ manifold. The
single-photon Rabi frequency of leg $i$ on the transition to an excited state
$e$ is $\Omega_{i,e}=d_e E_i/\hbar$, and with $I_i=\tfrac12\varepsilon_0cE_i^2$
\begin{equation}
\Omega_{i,e}^2=\frac{2\,d_e^2}{\varepsilon_0 c\,\hbar^2}\,I_i
\equiv d_e^2\,\kappa\,I_i,\qquad
\kappa=\frac{2}{\varepsilon_0 c\,\hbar^2} .
\end{equation}
For $|\Delta_e|\gg\Omega,\Gamma$ the excited states are eliminated
adiabatically. Each of them contributes one Raman path, and the paths add
coherently:
\begin{equation}
\Omega=\sum_e\frac{\Omega_{1,e}\,\Omega_{2,e}}{2\Delta_e}
=\underbrace{\kappa\sqrt{\beta(1-\beta)}
\sum_{e}\frac{d_{2e}\,d_{3e}}{2\Delta_e}}_{\textstyle C_\mathrm{rabi}}\;I ,
\end{equation}
The index $e$ labels an \emph{intermediate state}, i.e.\ a hyperfine level of
the $5P$ manifold,
\begin{equation}
|e\rangle=|5P_{J'},F',m'\rangle,\qquad
J'=\tfrac12\ (D_1,\ F'=2,3),\quad J'=\tfrac32\ (D_2,\ F'=1,2,3,4),
\end{equation}
so six terms for $\sigma^+$ light, where $m'=m+1$ is fixed by the polarisation.
$d_{2e}=\langle e|\hat d_{q}|F{=}2,m\rangle$ and
$d_{3e}=\langle e|\hat d_{q}|F{=}3,m\rangle$ are the \emph{signed} dipole matrix
elements, and
\begin{equation}
\Delta_e=\Delta+\left(\nu_{D_1}-\nu_\mathrm{line}\right)-\Delta_\mathrm{hfs}(F')
\end{equation}
is the detuning of the two-photon resonant pair from that level, measured from
the $5P_{1/2}$ centroid. The signs matter: the $F'$ contributions partly cancel,
and dropping them is wrong by tens of percent as soon as $|\Delta|$ approaches
the excited-state hyperfine splitting.

The formula above is evaluated with atomic data from the ARC package: the
signed matrix elements come from \texttt{getDipoleMatrixElementHFS}, the
hyperfine shifts from \texttt{getHFSCoefficients}/\texttt{getHFSEnergyShift},
the line centroids from \texttt{getTransitionFrequency} and the linewidths from
\texttt{getStateLifetime}. Only the elimination itself is coded by hand, in the method \texttt{coeff()}
of the module \texttt{rb85\_raman}. The same elimination gives the light shift of
each ground state and the photon scattering rate,
\begin{equation}
\delta_F=\kappa\sum_e d_{Fe}^2
\left[\frac{\beta}{4\Delta_e^{(1)}}+\frac{1-\beta}{4\Delta_e^{(2)}}\right]I,
\qquad
\Gamma_\mathrm{sc}=\kappa\sum_e\Gamma_e\,d_{Fe}^2
\left[\frac{\beta}{4(\Delta_e^{(1)})^2}+\dots\right]I ,
\end{equation}
so that all three are linear in $I$:
\begin{equation}
\Omega=C_\mathrm{rabi}I,\qquad
\delta_\mathrm{LS}=\delta_3-\delta_2\equiv C_\mathrm{shift}I,\qquad
\Gamma_\mathrm{sc}=C_\mathrm{scatter}I .
\end{equation}
Their ratio $\eta=\delta_\mathrm{LS}/\Omega$ is therefore a pure number,
independent of intensity, position and time. For $\sigma^+/\sigma^+$
co-propagating light on the clock states of Rb-85 and $\beta=\tfrac12$ one finds
$C_\mathrm{rabi}=11.66\,\mathrm{rad\,s^{-1}}/(\mathrm{W/m^2})$ and
$\eta=-0.061$ at $\Delta=+50$\,GHz, and $71.98$ and $-0.442$ at
$\Delta=-8$\,GHz. To leading order
$\eta\simeq-\omega_\mathrm{HF}/\Delta$, since the differential shift is the
difference of two denominators separated by the ground hyperfine splitting.

\paragraph{Step 3: the pulse area, without any reference.}
A rectangular pulse of length $T_p$ starting at $t_0$ accumulates the area
$\theta=\int\Omega\,\mathrm dt$. Three substitutions turn this into the working
formula. First the atomic law of Step 2, then the scale of Step 1, and finally
$s$ itself:
\begin{align}
\theta(\mathbf r,t_0)
&=\int_{t_0}^{t_0+T_p}\!\!\Omega(\mathbf r,t)\,\mathrm dt
&&\text{definition}\\
&=C_\mathrm{rabi}\int_{t_0}^{t_0+T_p}\!\! I(\mathbf r,t)\,\mathrm dt
&&\Omega=C_\mathrm{rabi}I\\
&=C_\mathrm{rabi}\,s\int_{t_0}^{t_0+T_p}\!\!\tilde I(\mathbf r,t)\,\mathrm dt
&&I=s\,\tilde I\\
&=C_\mathrm{rabi}\,\frac{P}{\tilde P_\mathrm{prof}}
  \int_{t_0}^{t_0+T_p}\!\!\tilde I(\mathbf r,t)\,\mathrm dt
&&s=P/\tilde P_\mathrm{prof} .
\end{align}
With $\tilde P_\mathrm{prof}=\int\langle\tilde I\rangle_t\,\mathrm d^2r$ this is
\begin{equation}
\boxed{\;
\theta(\mathbf r,t_0)
=C_\mathrm{rabi}\,P\;
\frac{\displaystyle\int_{t_0}^{t_0+T_p}\!\tilde I(\mathbf r,t)\,\mathrm dt}
     {\displaystyle\int\!\langle\tilde I\rangle_t\,\mathrm d^2r}\; }
\end{equation}
This is the whole physics. Only three things enter: the atomic coefficient, the
power, and the shape of the pattern, the latter as the ratio of a time integral
to an area integral. The arbitrary mode normalisation cancels in that ratio,
and no region, no $A_\mathrm{eff}$ and no $f_\mathrm{Rabi}$ appear. Units:
$[C_\mathrm{rabi}]\,\mathrm{W}\cdot\mathrm{s}/\mathrm{m^2}=\mathrm{rad}$.

\paragraph{Why $f_\mathrm{Rabi}$ may still be used.}
It is not an extra assumption but a name for the scale. Pick any reference
intensity $I_\mathrm{ref}$ -- in this work the time average over the evaluation
region -- and define
\begin{equation}
A_\mathrm{eff}=\frac{\int\langle I\rangle_t\,\mathrm d^2r}{I_\mathrm{ref}}
=\frac{\tilde P_\mathrm{prof}}{\langle\tilde I\rangle_\mathrm{ref}},
\qquad
I_\mathrm{ref}=\frac{P}{A_\mathrm{eff}},
\qquad
2\pi f_\mathrm{Rabi}\equiv C_\mathrm{rabi}I_\mathrm{ref} .
\end{equation}
Substituting $P=I_\mathrm{ref}A_\mathrm{eff}$ into the boxed equation gives
\begin{equation}
\theta(\mathbf r,t_0)
=2\pi f_\mathrm{Rabi}\,
\frac{1}{\langle\tilde I\rangle_\mathrm{ref}}
\int_{t_0}^{t_0+T_p}\!\tilde I(\mathbf r,t)\,\mathrm dt ,
\end{equation}
which is the same number: $A_\mathrm{eff}$ cancels against
$\tilde P_\mathrm{prof}$. Choosing a different region changes $I_\mathrm{ref}$,
$A_\mathrm{eff}$ and $f_\mathrm{Rabi}$ together and leaves $\theta$ untouched
(checked numerically to $10^{-14}$). $f_\mathrm{Rabi}$ is the Rabi frequency one
would measure by calibrating on the time averaged profile, without a trigger,
and $\Omega(\mathbf r,t)=C_\mathrm{rabi}I(\mathbf r,t)$ is what the atom
actually sees -- a quantity that varies in space and, through the beating,
strongly in time.

\paragraph{Where the sinc comes from.}
The time integral in the box is the only part that still depends on the
beating, and it can be done term by term. Inserting the Fourier series
$\tilde I=\sum_d D_d(\mathbf r)\,e^{i2\pi df_0t}$, each order contributes
\begin{equation}
\int_{t_0}^{t_0+T_p}\!\! e^{i2\pi df_0t}\,\mathrm dt
=\frac{e^{i2\pi df_0(t_0+T_p)}-e^{i2\pi df_0t_0}}{i\,2\pi df_0} .
\end{equation}
Pulling out the phase of the window \emph{centre}, $e^{i2\pi
df_0(t_0+T_p/2)}$, makes the bracket symmetric,
\begin{equation}
=e^{i2\pi df_0\left(t_0+T_p/2\right)}\,
\frac{e^{i\pi df_0T_p}-e^{-i\pi df_0T_p}}{i\,2\pi df_0}
=e^{i2\pi df_0\left(t_0+T_p/2\right)}\,
\frac{\sin\!\left(\pi df_0T_p\right)}{\pi df_0} ,
\end{equation}
where $e^{ix}-e^{-ix}=2i\sin x$ was used. Writing the last factor as $T_p$
times $\operatorname{sinc}(z)=\sin(\pi z)/(\pi z)$ gives
\begin{equation}
\int_{t_0}^{t_0+T_p}\!\tilde I\,\mathrm dt
=T_p\sum_d D_d(\mathbf r)\,\operatorname{sinc}(d f_0T_p)\,
e^{i2\pi df_0\left(t_0+T_p/2\right)} .
\end{equation}
The three factors read directly: $T_p$ is the length of the window, the
exponential says that only the \emph{centre} of the window matters for the
phase, and the sinc is the Fourier transform of the window itself, evaluated at
the beat frequency $df_0$. Integrating over a window of length $T_p$ is a
convolution with a boxcar, and in frequency space a convolution is a
multiplication -- with the transform of the box. The $d=0$ term is the limit
$\operatorname{sinc}(0)=1$: the time average always survives, whatever the
pulse length. Because $D_{-d}=D_d^*$ the sum is real,
\begin{equation}
\begin{split}
\int_{t_0}^{t_0+T_p}\!\tilde I\,\mathrm dt
=T_p\Bigl[\,&D_0(\mathbf r)\\
&+2\sum_{d>0}\left|D_d(\mathbf r)\right|
\operatorname{sinc}(df_0T_p)\,
\cos\!\left(2\pi df_0\left(t_0+\tfrac{T_p}{2}\right)+\arg D_d\right)\Bigr].
\end{split}
\end{equation}
This is the same expression as a camera exposure of duration $T_p$; a pulse and
an exposure differ only in their length. Two limits:
\begin{itemize}
\item $T_p\gg T_0$: every $d\neq0$ is damped away and
      $\theta\to C_\mathrm{rabi}\,s\,\langle\tilde I\rangle_t\,T_p$. The pulse
      sees the time average.
\item $T_p\ll T_0$: $\operatorname{sinc}\to1$. The pulse does not average the
      beating, it \emph{samples} it, and $\theta$ depends on the trigger
      instant $t_0$ -- on the time scale of the fast beats, not of $T_0$.
\end{itemize}

\paragraph{The pulse area in closed form.}
Inserting the boxcar result into the boxed equation gives the expression the
code evaluates,
\begin{equation}
\theta(\mathbf r,t_0)
=\frac{C_\mathrm{rabi}\,P\,T_p}{\tilde P_\mathrm{prof}}
\sum_d D_d(\mathbf r)\,\operatorname{sinc}(df_0T_p)\,
e^{i2\pi df_0\left(t_0+T_p/2\right)} ,
\end{equation}
or, using $D_{-d}=D_d^*$, in manifestly real form
\begin{equation}
\begin{split}
\theta(\mathbf r,t_0)=\frac{C_\mathrm{rabi}PT_p}{\tilde P_\mathrm{prof}}
\Bigl[\,&D_0(\mathbf r)\\
&+2\sum_{d>0}\left|D_d(\mathbf r)\right|\operatorname{sinc}(df_0T_p)\,\times\\
&\qquad\cos\!\left(2\pi df_0\!\left(t_0+\tfrac{T_p}{2}\right)+\arg D_d\right)
\Bigr].
\end{split}
\end{equation}
The first term is the time average and carries no trigger dependence; every
other order is damped by its sinc and swings with the trigger. A $\pi$ pulse
requires $\theta=\pi$, so the power that achieves it at a given trigger is
$P_\pi=P\,\pi/\theta$.

\paragraph{Averages over the atom are integrals.}
$\theta(\mathbf r)$ is a field. What the atom experiences is the average over
its position probability density, a proper integral,
\begin{gather}
\langle\theta\rangle_W=\int\! W(\mathbf r)\,\theta(\mathbf r)\,\mathrm d^2r,
\qquad \int\! W(\mathbf r)\,\mathrm d^2r=1,\\
W(\mathbf r)=\frac{1}{2\pi\sigma^2}
\exp\!\left[-\frac{(\mathbf r-\mathbf r_0)^2}{2\sigma^2}\right],
\end{gather}
with the thermal width
$\sigma^2=\frac{\hbar}{2m\omega_r}\coth\frac{\hbar\omega_r}{2k_BT}$
($\sigma=108$\,nm at $\nu_r=60.4$\,kHz and $T=17\,\mu$K). A hard region is the
special case $W=1/A_R$ inside and $0$ outside. On the uniform computation grid
the integral becomes the weighted sum
$\sum_i w_i\theta_i$ with $w_i=W_i/\sum_iW_i$, because the cell area cancels
between numerator and denominator; that is how the code evaluates it, on a
dedicated fine grid of $\pm4\sigma$, since the global grid is far too coarse for
$\sigma\approx0.1\,\mu$m.

Because $\theta$ is \emph{linear} in the coefficients $D_d$, the spatial average
can be pulled inside the sum. With the atom weighted coefficients
\begin{equation}
c_d=\int\! W(\mathbf r)\,D_d(\mathbf r)\,\mathrm d^2r
\qquad\text{(one complex number per beat order)}
\end{equation}
the mean area becomes
\begin{equation}
\langle\theta\rangle_W(t_0)
=\frac{C_\mathrm{rabi}PT_p}{\tilde P_\mathrm{prof}}
\sum_d c_d\,\operatorname{sinc}(df_0T_p)\,e^{i2\pi df_0(t_0+T_p/2)} .
\end{equation}
The $c_d$ are computed once, and the whole scan over the trigger $t_0$ is then a
single matrix-vector product instead of one field evaluation per delay. The
spread of the area over the atom needs the map itself,
\begin{equation}
\sigma_\theta^2=\int\! W\,\theta^2\,\mathrm d^2r-\langle\theta\rangle_W^2 ,
\qquad
U(\theta)=\frac{\sigma_\theta}{\langle\theta\rangle_W} ,
\end{equation}
and the excitation must be averaged \emph{after} the non-linear step,
\begin{equation}
\langle p_\uparrow\rangle_W=\int\! W(\mathbf r)\,
\sin^2\!\frac{\theta(\mathbf r)}{2}\,\mathrm d^2r
\;\neq\;\sin^2\!\frac{\langle\theta\rangle_W}{2} .
\end{equation}

\paragraph{What the atom does with the area.}
For a resonant two-level system with a real $\Omega(t)$ the Hamiltonian
commutes with itself at different times, so the excitation depends on the
accumulated area alone -- exactly, not approximately:
\begin{equation}
p_\uparrow(\mathbf r)=\sin^2\frac{\theta(\mathbf r)}{2},
\qquad
p_\uparrow(\mathbf r)=\frac{1}{1+\eta^2}
\sin^2\!\left(\sqrt{1+\eta^2}\,\frac{\theta(\mathbf r)}{2}\right)
\ \text{with the light shift},
\end{equation}
and the measured signal is $\int W p_\uparrow\,\mathrm d^2r$: $\sin^2$ is
non-linear, so the average is taken \emph{after} it. The spread of the area
itself, $U(\theta)=\sigma_W(\theta)/\langle\theta\rangle_W$, is what limits the
fidelity. Scattering multiplies the signal by $e^{-\Gamma_\mathrm{sc}T_p}$.

\paragraph{Working point.}
$3\times4$, Airy, waist $1.04\,\mu$m, $r_x/r_y=0.97/1.16$, all tone phases
zero, $\Delta=+50$\,GHz, $P=10\,\mu$W in the profile, $T_p=1\,\mu$s, atom
weighted ($\sigma=108$\,nm):

\begin{center}
\begin{tabular}{lr}
\hline
$\tilde P_\mathrm{prof}$ (simulation units) & $2.615\times10^{-11}$\,m$^2$\\
$\langle\tilde I\rangle_W$ (simulation units) & $1.913$\\
effective area $A_\mathrm{eff}$ & $13.67\,\mu$m$^2$\\
reference intensity $I_\mathrm{ref}=P/A_\mathrm{eff}$ & $73.2$\,W/cm$^2$\\
nominal Rabi frequency $f_\mathrm{Rabi}$ & $1.358$\,MHz\\
area at the reference, $2f_\mathrm{Rabi}T_p$ & $2.72\,\pi$\\
area at the best trigger, $\langle\theta\rangle_W$ & $11.54\,\pi$\\
mean Rabi frequency during that pulse & $5.77$\,MHz\\
$\sigma_\theta/\langle\theta\rangle_W$ there & $0.72\,\%$\\
power for a $\pi$ pulse of $1\,\mu$s at that trigger & $0.87\,\mu$W\\
\hline
\end{tabular}
\end{center}

Both routes -- the boxed formula and the detour over $f_\mathrm{Rabi}$ -- give
the same $11.539\,\pi$. The factor $4.25$ between the triggered and the
reference area is the rephasing peak of the twelve tones at phases zero.

\paragraph{Not included.}
Static two-photon detuning and magnetic field shifts, any uncompensated
gradient of the light shift, spontaneous emission during the pulse, the Zeeman
substructure, atomic motion during the pulse, finite pulse edges and the
acoustic fill time of the AOD, polarisation and vector light shifts,
intermodulation in the AOD and its diffraction efficiency. $P$ is the power
\emph{in the profile}, i.e. after both AODs.
